\section{Introduction}

\subsection{Motivation and Significance}

Computing an exact Nash equilibrium of a bimatrix game is \PPAD-complete \cite{DGP09}, making efficient approximation a central problem in algorithmic game theory. Deligkas, Fasoulakis, and Markakis advance the polynomial-time guarantee for the stronger notion of well-supported Nash equilibrium from $0.6528$ to $1/2+\delta$ for every fixed $\delta>0$ \cite{12WSNE}. Beyond the numerical improvement, the paper identifies a comparatively simple structural route to the result: zero-sum equilibria separate the instance into cases handled by linear feasibility or constant-support enumeration. Its payoff-query extension shows that essentially the same threshold remains attainable with near-linear query complexity for fixed approximation parameters \cite[Section~4]{12WSNE}.

We selected this paper mainly for its landmark significance in the game theory. Besides, its appearance in the lecture notes stimulated our interest. Another reason is that its relatively concise presentation made the central argument seemingly more accessible, though we later realized that this seemed to be because it had skipped steps.

\subsection{Our Contribution}

Our contribution is expository and corrective. For both Theorem~\ref{thm:finding-bimat-game-12_del-wsne-is-poly-time} and its query-efficient counterpart, Theorem~\ref{thm:query-efficient-version}, we reconstruct deductions that are partly omitted or compressed in the source. We display the algorithms in a more detailed manner, elaborating on each step. In particular, we make explicit how the zero-sum maximin guarantees, the linear feasibility systems, and the constant-support enumeration establish correctness in every branch. We also supply the detailed  case-analysis arguments collected in Appendix~\ref{app:case-analysis-details}.

We further audit the randomized implementation and identify \textbf{probability-accounting errors} in the source. The stated success probability covers only one invocation of the zero-sum routine, although Algorithm~2 invokes it twice and additionally samples a small-support profile. Besides, we take sampling failures into consideration. Incorporating a sampling failure budget $\eta$, we derive the corrected bound $p_n^2(1-\eta)$ and re-establish the sampling, query, and running-time bounds, with the paper's principal asymptotic conclusions unchanged but more rigorous. We further uncover a query-accounting gap: the branch tests of Algorithm~2 require exact evaluations of $\mathcal{R}(\mathbf{x}^*,\mathbf{y}^*)$ and $\mathcal{C}(\hat{\mathbf{x}},\hat{\mathbf{y}})$, which the zero-sum routine does not return and which may cost $\Theta(n^2)$ queries. We repair this by substituting statistical estimates with explicit slack $\gamma\le\varepsilon$, preserving the asymptotic query bound and the approximation guarantee at the price of a reduced success probability. We found this repair independently, with the assistance of Codex; we later learned that an earlier preprint by Zhengyang Liu \cite{liu2026supportcontractionwellsupportednash} repairs the same gap via a stronger support-contraction method, improving the query count to $O(n\log n/\varepsilon^2)$. Section~\ref{sec:source-corrections} also records several minor errors.

